% 本文件是中期考核报告的填写入口。
% 正式使用时，优先替换下面 \midtermsetup 中的示例值。
% 正文内容较长时，默认 \printmidterm 会自动分页，不需要手动控制 6 页。

\midtermsetup{
  info = {
    student-id = {2023123456},
    name = {张三},
    first-level-discipline = {电子科学与技术},
    second-level-discipline = {电路与系统},
    supervisor = {李四 教授},
    school = {电子工程学院},
    report-date = {2026年7月},
    title = {基于深度学习的目标检测\\方法研究}
  },
  content = {
    source = {导师科研项目},
    research-goal = {
      本研究面向若干实际应用场景，围绕关键问题开展理论分析、
      算法设计与实验验证。
    },
    completed-work = {
      目前已完成相关文献调研、基础模型复现和初步实验分析。
    },
    deviation = {
      目前研究内容与开题报告基本一致。
    },
    next-plan = {
      后续将继续完善实验方案，补充对比实验，并完成论文主要章节撰写。
    },
    achievements = {无。}
  },
  review = {
    supervisor-comment = {
      该生按计划推进课题研究，已完成阶段性工作，同意进入下一阶段。
    },
    supervisor-sign = {李四},
    supervisor-year = {2026},
    supervisor-month = {7},
    supervisor-day = {3},
    committee-comment = {
      中期考核通过，建议根据专家意见继续完善研究内容和论文写作。
    },
    chair-sign = {王五},
    member-sign = {赵六、钱七},
    committee-year = {2026},
    committee-month = {7},
    committee-day = {3}
  }
}

% 旧接口仍可用，例如：
% \renewcommand{\XueHao}{2023123456}
% \renewcommand{\XingMing}{张三}
